\chapter{Behaviorální stromy}

Tato kapitola popisuje základní principy behaviorálních stromů a jejich využití při návrhu chování NPC. 
Nejprve je vysvětleno, k čemu behaviorální stromy slouží a jakým způsobem se vykonávají. 
Poté jsou popsány základní typy uzlů, ze kterých se skládají. 
Závěr kapitoly shrnuje hlavní výhody a omezení tohoto přístupu s ohledem na návrh chování postav v UE5.

\section{Úloha behaviorálních stromů}

Behaviorální strom je způsob zápisu rozhodovací logiky autonomního agenta. 
V oblasti počítačových her se tímto agentem obvykle rozumí NPC, která na základě aktuální situace vybírá vhodné chování. 
Může například procházet herní prostředí, reagovat na hráče, komunikovat s jinou postavou nebo provádět interakce s objekty v herním prostředí.

Základní myšlenkou behaviorálního stromu je rozdělit chování na menší části. 
Místo jedné rozsáhlé podmínkové struktury je rozhodování popsáno pomocí uzlů uspořádaných do stromu. 
Některé uzly určují pořadí vyhodnocování, jiné ověřují podmínky a další vykonávají konkrétní akce. 
Díky tomu je možné složitější chování sestavit z jednodušších a lépe oddělených částí.\cite[s.~6]{colledanchise2018}

Tento způsob návrhu je vhodný zejména pro herní umělou inteligenci, protože chování postav musí být často reaktivní. 
Postava se nemůže rozhodnout pouze jednou a poté pokračovat bez ohledu na změny okolí. 
Naopak musí průběžně reagovat na pohyb hráče, změnu stavu objektů, dostupnost cesty nebo výsledek právě prováděné akce. 
Behaviorální strom umožňuje takové změny zachytit opakovaným vyhodnocováním rozhodovací logiky.\cite[s.~5]{colledanchise2018}

V této práci jsou behaviorální stromy důležité především proto, že představují hlavní prostředek pro řízení chování NPC v UE5. 
Pomocí nich lze jednotlivé typy chování oddělit do samostatných větví a následně určit, za jakých podmínek se mají spouštět.

\section{Základní princip vykonávání}

Behaviorální strom se vykonává od kořenového uzlu. 
Z něj je v pravidelných intervalech vysílán signál označovaný jako \textit{tick}. 
Tento signál postupuje stromem podle pravidel jednotlivých uzlů. 
Uzel je tedy vyhodnocen pouze tehdy, pokud k němu \textit{tick} dorazí.\cite[s.~6]{colledanchise2018}

Každý uzel po vyhodnocení vrací svému rodiči výsledek. 
V běžném pojetí behaviorálních stromů se používají tři základní návratové stavy. 
Stav \textit{Success} znamená, že uzel byl úspěšně dokončen. Stav \textit{Failure} vyjadřuje, že uzel nemohl svůj úkol splnit. 
Stav \textit{Running} se používá v případě, kdy akce stále probíhá a její výsledek zatím není známý.\cite[s.~9, tab.~1.1]{colledanchise2018}

Tyto návratové stavy tvoří základ komunikace mezi jednotlivými částmi stromu. 
Například akce přesunu na určité místo většinou nemůže skončit během jediného snímku hry. 
Dokud se postava pohybuje k cíli, uzel může vracet stav \textit{Running}. 
Jakmile je cíl dosažen, vrátí \textit{Success}. 
Pokud se cestu nepodaří najít, výsledek bude \textit{Failure}.

Opakované vyhodnocování stromu umožňuje postavě reagovat na aktuální stav hry. 
Pokud se například během pohybu změní podmínky v okolí, může být při dalším vyhodnocení zvolena jiná větev. 
Tím se behaviorální strom liší od jednoduchého skriptu, který pouze provede předem danou posloupnost příkazů bez průběžného přehodnocování situace.\cite[s.~9]{colledanchise2018}

\section{Struktura behaviorálního stromu}
Behaviorální strom je tvořen uzly, které jsou mezi sebou propojeny vztahem rodič--potomek. 
Výjimkou je kořenový uzel, který má vazby pouze na své potomky. 
Kořenový uzel stojí na začátku vyhodnocování a jednotlivé větve stromu popisují možné části chování.\cite[s.~6, kap.~1.3]{colledanchise2018}

Z pohledu struktury se uzly dělí do tří kategorií: kompozitní, dekorátorové a listové. 
Kompozitní uzly řídí tok vykonávání a určují, v jakém pořadí a za jakých podmínek se budou potomci vyhodnocovat. 
Dekorátorové uzly mají právě jednoho potomka a modifikují jeho chování nebo návratový stav. 
Listové uzly se nacházejí na koncích větví a představují buď podmínkové ověření stavu, nebo konkrétní akci. 
Podrobný popis každé kategorie je uveden v sekci \ref{sec:behavior_tree_node_types}.

Kompozitní uzly neurčují konkrétní akci ve světě hry, ale rozhodují o tom, která část stromu bude vyhodnocena. 
Typicky určují pořadí potomků nebo vybírají první větev, jejíž podmínky jsou splněny. 
Díky nim lze ve stromu vyjádřit například prioritu jednotlivých chování.

Podmínkové listové uzly slouží k ověření určitého stavu. 
V herním prostředí může jít například o kontrolu, zda postava vidí hráče, zda je v dosahu objekt, nebo jestli existuje platný cíl pohybu. 

Akční listové uzly představují samotné činnosti postavy. 
Mohou spustit pohyb, vybrat nový cíl, změnit hodnotu v \textit{Blackboard}, přehrát animaci nebo provést interakci s objektem. 
V UE5 jsou pro tyto části používány především uzly typu \textit{Task}.\cite[s.~76]{butler2023}

\section{Základní typy uzlů}
\label{sec:behavior_tree_node_types}

\subsection{Kompozitní uzly}

\textit{Kompozitní} uzly jsou řídicí uzly, které mají jednoho nebo více potomků. 
Jejich úlohou je určit, v jakém pořadí a za jakých podmínek se budou potomci vyhodnocovat. 
Mezi nejčastěji používané \textit{Kompozitní} uzly patří \textit{Sequence} a \textit{Selector} (v literatuře též zvaný jako \textit{Fallback}). 
V některých systémech se používá také uzel \textit{Parallel}.\cite[s.~6--7]{colledanchise2018}

\subsubsection{Sequence}

Uzel \textit{Sequence} slouží k postupnému vykonání několika kroků. 
Vyhodnocuje své potomky zleva doprava a pokračuje pouze tehdy, pokud předchozí potomek skončil úspěchem. 
Pokud některý potomek vrátí stav \textit{Failure}, selže celá sekvence. 
Pokud potomek vrátí \textit{Running}, znamená to, že daný krok ještě probíhá, a stejný stav vrací také celá sekvence.
Uzel \textit{Sequence} vrací návratovou hodnotu \textit{Success} pouze v tom případě, když všichni potomci tohoto uzlu mají návratovou hodnotu \textit{Success}.~\cite[s.~6]{colledanchise2018}

Tento uzel je vhodný pro chování, u kterého musí být splněno více kroků v přesném pořadí. 
Příkladem může být interakce NPC s předmětem. 
Postava nejprve ověří, zda je předmět dostupný, poté se k němu přesune a nakonec provede samotnou akci. 
Pokud už první podmínka není splněna, nemá smysl pokračovat v dalších krocích.

\subsubsection{Selector}

Uzel \textit{Selector} slouží k výběru první vhodné možnosti. 
Stejně jako \textit{Sequence} prochází své potomky zleva doprava, ale používá opačnou logiku. 
Jakmile některý potomek uspěje nebo stále probíhá, \textit{Selector} se zastaví a vrátí odpovídající stav. 
Selže pouze tehdy, když selžou všichni jeho potomci.~\cite[s.~7]{colledanchise2018}

\textit{Selector} se často používá pro vyjádření priorit. 
První větev může například popisovat reakci na hráče, druhá komunikaci s jinou postavou a třetí běžné pohybování po scéně. 
Strom se nejdříve pokusí vykonat chování s nejvyšší prioritou. 
Pokud pro něj nejsou splněny podmínky, pokračuje k další možnosti.

\subsubsection{Parallel}

Uzel \textit{Parallel} umožňuje vyhodnocovat více potomků současně.
Používá se v případech, kdy má agent provádět více činností najednou nebo kdy má během jedné akce průběžně sledovat další podmínku.
Chování uzlu je řízeno dvěma parametry: $N$ označuje celkový počet přímých potomků a $M$ je konfigurovatelný práh úspěchu ($1 \leq M \leq N$). 
Uzel vrátí stav \textit{Success}, pokud alespoň $M$ potomků uspěje, a stav \textit{Failure}, pokud více než $N - M$ potomků selže dříve, než je práh $M$ dosažen.\cite[s.~7]{colledanchise2018}

V UE5 existuje uzel \textit{Simple Parallel}. 
Ten umožňuje spustit hlavní úlohu a vedle ní vyhodnocovat další větev stromu.\cite[s.~76]{butler2023} 
Prakticky může jít například o~situaci, kdy se postava přesouvá k cíli a zároveň kontroluje, zda se nezměnil stav okolí.

\subsection{Dekorátorové uzly}
\label{subsec:decorators}

Dekorátorové uzly představují speciální typ řídicích uzlů, které mají právě jednoho potomka a slouží k úpravě jeho chování nebo návratového stavu.\cite[s.~8--9]{colledanchise2018}
Na rozdíl od ostatních řídicích uzlů, které určují tok vykonávání mezi více potomky, dekorátory modifikují výsledek nebo podmínky vykonání jediné podřízené větve.

V obecné teorii behaviorálních stromů mohou dekorátory například měnit návratové hodnoty uzlů (např. převrácení výsledku \textit{Success} na \textit{Failure}) 
nebo ovlivňovat způsob vykonávání, například omezením počtu opakování určité akce.\cite[s.~9]{colledanchise2018}

V prostředí UE5 mají dekorátory poněkud odlišnou vizuální strukturu. 
Jsou obvykle připojeny ke konkrétní větvi stromu a slouží především jako podmínky určující, zda může být daná větev vykonána.\cite[s.~93]{butler2023}
Nemají tedy vlastní samostatný vizuální uzel, jako tomu je u ostatních uzlů.
Typickým příkladem je kontrola hodnoty uložené v \textit{Blackboard}. 
Pokud je například dostupná reference na hráče, může se vykonat větev popisující reakci na hráče. 
V opačném případě je tato větev přeskočena a strom pokračuje jinou alternativou.

\subsection{Listové uzly}

Listové uzly se nacházejí na koncích větví behaviorálního stromu a nemají žádné potomky. 
Jejich úkolem je buď ověřit určitý stav, nebo provést konkrétní akci.\cite[s.~6]{colledanchise2018}

Podmínkové uzly kontrolují stav prostředí nebo interních dat postavy. 
Mohou například ověřit, zda je hráč v dosahu, zda je dostupný objekt pro interakci nebo zda postava již dokončila předchozí úkol. 
Výsledkem takového uzlu je zpravidla stav \textit{Success} nebo \textit{Failure}.\cite[s.~8]{colledanchise2018}
Stav \textit{Running} se u podmínkového uzlu nevyskytuje kvůli rozhodovací podstatě tohoto uzlu.

Akční uzly vykonávají konkrétní činnost. 
V UE5 se jedná zejména o uzly typu \textit{Task}. 
Ty mohou například spustit přesun pomocí navigačního systému, vybrat náhodný bod ve scéně, změnit hodnotu v \textit{Blackboard} 
nebo zavolat funkci implementovanou v \textit{Blueprint} či C++ třídě. 
Právě akční uzly propojují rozhodovací logiku behaviorálního stromu se skutečným chováním postavy ve hře.

Přehled návratových stavů všech popsaných typů uzlů je uveden v~tabulce~\ref{tab:BehaviorTreeNodeTypes}.

\begin{table}[H]
\centering
\small
\setlength{\fboxsep}{3pt}
\setlength{\fboxrule}{0.5pt}

\begin{tabularx}{\textwidth}{l c X X X}
\hline
\textbf{Typ uzlu} & \textbf{Symbol} & \textbf{Success} & \textbf{Failure} & \textbf{Running} \\
\hline
Selector & \fbox{?} & Pokud alespoň jeden potomek uspěje & Pokud všichni potomci selžou & Pokud alespoň jeden potomek vrátí Running \\
\hline
Sequence & \fbox{$\rightarrow$} & Pokud všichni potomci uspějí & Pokud alespoň jeden potomek selže & Pokud alespoň jeden potomek vrátí Running \\
\hline
Parallel & \fbox{$\rightrightarrows$} & Pokud $\geq M$~potomků uspěje & Pokud $> N - M$~potomků selže & V ostatních případech \\
\hline
Action & \fbox{text} & Po dokončení & Pokud nelze dokončit & Během provádění \\
\hline
Condition & \btboxround{text} & Pokud platí (true) & Pokud neplatí (false) & Nikdy \\
\hline
Decorator & $\diamond$ & Vlastní definice & Vlastní definice & Vlastní definice \\
\hline
\end{tabularx}

\caption{Typy uzlů v behaviorálním stromě.
Přeloženo podle~\cite[s.~9, tab.~1.1]{colledanchise2018}.}
\label{tab:BehaviorTreeNodeTypes}
\end{table}

Výše uvedená tabulka popisuje, jak fungují návratové hodnoty jednotlivých typů uzlů, ale neméně důležité je jejich vzájemné uspořádání ve stromě. 
Následující diagram ukazuje, jak lze popsané uzly kombinovat na konkrétním příkladu.

\begin{figure}[H]
\centering
\begin{forest}
  for tree={
    draw, rectangle, align=center, text centered,
    minimum width=3cm, minimum height=0.8cm,
    text width=2.8cm, font=\small,
    s sep=3mm, l sep=20mm,
    edge={-stealth}, anchor=north, inner sep=4pt,
  }
  [{ROOT}, fill=gray!30
    [{Selector~(?)}, fill=gray!15
      [{Sequence~($\rightarrow$)}, fill=gray!10
        [{Je hráč\\viditelný?},
            draw, rounded corners=6pt, align=center,
            text width=2.4cm, fill=orange!20]
        [{Uprav chování\\dle postupu\\hráče}]
      ]
      [{Sequence~($\rightarrow$)}, fill=gray!10
        [{Je v dosahu\\jiné NPC?},
            draw, rounded corners=6pt, align=center,
            text width=2.4cm, fill=orange!20]
        [{Zahaj komunikaci\\s~NPC}]
      ]
      [{Přesuň\\předmět}]
    ]
  ]
\end{forest}
\caption{Zjednodušená struktura behaviorálního stromu pro NPC.}
\label{fig:bt_npc_example}
\end{figure}

Obrázek~\ref{fig:bt_npc_example} znázorňuje, jak lze tři navrhovaná chování přirozeně uspořádat do jediného stromu. 
Uzel \textit{Selector} prochází větve zleva doprava a zkouší vždy chování s nejvyšší prioritou. 
Nejprve se pokusí o adaptaci vůči hráči. 
Podmínka viditelnosti hráče musí být splněna, aby se \textit{Sequence} vykonala.
Pokud hráč není viditelný, strom přejde na komunikaci s jiným NPC.
Není-li ani to možné, vykoná se výchozí chování, kterým je přesun předmětu. 
Tato větev nemá podmínku, takže k ní Selector vždy dojde. 
Uzel může skončit úspěchem i selháním podle toho, jak skončí daná akce.


\section{Výhody behaviorálních stromů}

Jednou z hlavních výhod behaviorálních stromů je modularita. 
Chování lze rozdělit do menších částí, které mají jasně určenou odpovědnost. 
Jedna větev může řešit pohyb, jiná reakci na hráče a další například komunikaci mezi postavami. 
Pokud jsou tyto části navrženy obecně, lze je později použít i v jiném chování nebo u jiné postavy.\cite[s.~38]{colledanchise2018}

Další výhodou je přehlednost. 
Stromová struktura umožňuje sledovat, jakým způsobem se postava rozhoduje. 
To je užitečné zejména v UE5, kde lze behaviorální strom vytvářet a upravovat v grafickém editoru. 
Vývojář tak nemusí procházet rozsáhlý zdrojový kód, ale může pracovat s vizuální reprezentací jednotlivých větví.\cite[s.~38]{colledanchise2018}

Důležitá je také reaktivita. 
Protože se strom vyhodnocuje opakovaně, může postava reagovat na změny prostředí. 
Pokud se například během běžného pohybu objeví hráč, může strom při dalším vyhodnocení přejít do větve s vyšší prioritou. 
Díky tomu chování nepůsobí pouze jako pevně připravená posloupnost příkazů, ale může se přizpůsobovat aktuální situaci.\cite[s.~38]{colledanchise2018}

Behaviorální stromy také podporují postupný vývoj. 
Na začátku může být strom jednoduchý a obsahovat pouze několik základních větví. 
Později lze přidávat další podmínky, akce a reakce, aniž by bylo nutné přepsat celý systém rozhodování. 
To je výhodné při vývoji herní umělé inteligence, protože chování postav se často upravuje podle výsledků testování.\cite[s.~38]{colledanchise2018}

\section{Omezení behaviorálních stromů}

Behaviorální stromy mají také určitá omezení. 
U rozsáhlých stromů může postupně klesat přehlednost, zejména pokud nejsou jednotlivé větve vhodně pojmenované nebo pokud se v nich opakuje podobná logika. 
V takovém případě může být obtížné zjistit, proč postava zvolila právě určité chování.

Dalším problémem může být špatně navržené pořadí větví. 
Uzel \textit{Selector} upřednostňuje potomky podle jejich pořadí. 
Pokud je důležitá větev umístěna příliš nízko, strom se k ní nemusí dostat, protože dříve uspěje jiná možnost. 
Chyba tedy nemusí být v samotné akci, ale ve struktuře rozhodování.

Behaviorální strom také vyžaduje průběžné vyhodnocování podmínek. 
U menšího počtu postav to obvykle nepředstavuje problém, ale u většího množství agentů nebo složitějších podmínek může být nutné řešit výkon. 
Při návrhu stromu je proto vhodné zvážit, které kontroly je nutné provádět často a které lze vyhodnocovat méně pravidelně.\cite[s.~42]{colledanchise2018}

\section{Srovnání s jinými technikami implementace chování NPC}

Behaviorální stromy nejsou ani zdaleka jedinou metodou řízení chování autonomních agentů.
V této sekci budou představeny a porovnány behaviorální stromy se dvěma alternativními přístupy, které jsou používány v~herním vývoji.
Jedná se o konečné automaty (\textit{FSM}) a~systémy založenými na hodnotě užitku (\textit{Utility-based system}).


\subsection{Konečné automaty}

Konečný automat (Finite State Machine, zkráceně FSM) je historicky nejrozšířenější přístup k~implementaci chování NPC. 
Agent se v~každém okamžiku nachází v~jednom z~předem definovaných stavů a~přechází mezi nimi na základě podmínek, které jsou průběžně vyhodnocovány.~\cite[s.~318]{millington2006}. 
Pro jednoduché NPC, kterému stačí pár stavů pro jeho věrohodnou implementaci chování, FSM bohatě stačí.
Asi nejlepším příkladem takových NPC jsou obecně zvířata.

Hlavním omezením FSM je škálovatelnost kvůli tomu, že s~rostoucím počtem stavů roste i~počet nutných přechodů mezi nimi. 
Každý nový stav tak může vyžadovat definování přechodů ze všech stavů ostatních. 
Tento jev je označován jako problém exploze stavů~\cite[s.~5]{colledanchise2018}. 
Další nevýhodou je obtížná znovupoužitelnost, protože chování navržené pro jednoho agenta 
lze jen těžko přenést na jiného agenta s~odlišnou sadou stavů.

V~kontextu implementovaného systému by použití FSM vedlo k~výraznému nárůstu počtu přechodů mezi stavy, zejména při kombinaci percepce, hlídkování 
a~interakcí s~objekty. 
Takový model by byl obtížně rozšiřitelný a~méně přehledný při ladění chování.

Behaviorální stromy tento problém řeší hierarchickým uspořádáním chování do větví. 
Dílčí chování jsou izolována do samostatných podstromů, které lze sdílet mezi různými agenty nebo opakovaně skládat do komplexnějšího chování. 
Přidání nového chování nevyžaduje úpravu existujících přechodů, ale pouze přidání nové větve do stromu. 
Z~těchto důvodů jsou behaviorální stromy považovány za přirozenou náhradu FSM v~situacích, kdy počet chování přesáhne několik málo stavů~\cite[s.~5]{colledanchise2018}.

\subsection{Systémy založené na hodnotě užitku}

Systémy založené na hodnotě užitku (utility-based systems) přiřazují každé
dostupné akci číselné skóre vyjadřující její vhodnost v~dané situaci. 
Agent nejčastěji vykoná akci s~nejvyšším skóre~\cite[s.~379]{millington2006}. 
Hodnota akce se typicky počítá jako vážená kombinace více faktorů. 
Například vzdálenosti od hráče, úrovně zdraví agenta nebo dostupnosti předmětů v~prostředí.

Výhodou tohoto přístupu je plynulý a~přirozený výběr chování. Agent nepřepíná diskrétně mezi stavy, ale průběžně reaguje 
na jemné změny situace. Pro simulace, kde je žádoucí, aby NPC působila organicky a~ne mechanicky předvídatelně, může být tato vlastnost přínosem.

Nevýhodou je obtížnost návrhu a~ladění. Správné nastavení vah a~hodnotících funkcí pro každou kombinaci faktorů vyžaduje opakované testování. 
Výsledné chování nevyplývá z~přehledné struktury, ale z~numerického výpočtu, jehož průběh je těžší sledovat a~opravovat.

V~případě implementovaného systému by bylo nutné definovat hodnotící funkce pro kombinaci 
více faktorů (např. vzdálenost hráče, slyšitelnost nebo dostupnost interakcí), což by výrazně zkomplikovalo ladění chování
bez zřejmého přínosu oproti explicitní hierarchii behaviorálního stromu.

Oproti tomu behaviorální strom dává návrháři explicitní kontrolu nad prioritou chování. 
Strukturou větví přímo určuje, co má přednost a za jakých podmínek, bez nutnosti kalibrovat numerické váhy.

\subsection{Zhodnocení}

Oproti FSM poskytují behaviorální stromy lepší škálovatelnost a~přehlednost při rostoucí komplexitě chování. 
Ve srovnání se systémy založenými na hodnotě užitku umožňují explicitní řízení priorit a~jednodušší ladění bez nutnosti 
definování numerických hodnotících funkcí.

Zvolený přístup se ukázal jako vhodný zejména pro implementaci hierarchického chování NPC, kde jednotlivé 
stavy (\texttt{Passive}, \texttt{Investigating}, \texttt{Attacking}) a~jejich přechody 
přirozeně odpovídají struktuře behaviorálního stromu.